\section{Auto-évaluation et analyse réflexive}

\begin{guideline}
C'est la partie la plus importante et la plus personnelle du dossier. Elle démontre votre capacité à prendre du recul sur votre expérience. Soyez précis·e, honnête et utilisez des exemples concrets tirés de votre stage. Évitez les généralités.
\end{guideline}


\subsection{Auto-évaluation de la mission}

\begin{guideline}
Faites un bilan de votre mission : avez-vous atteint les objectifs fixés ? Quels résultats avez-vous produits ? Quelle valeur ajoutée avez-vous apportée (pour l'entreprise, l'équipe, les utilisateurs) ? Quantifiez quand c'est possible (ex. : gain de temps, taux d'adoption, réduction des erreurs...).
\end{guideline}

\instruction{[Dressez un bilan de votre mission : objectifs atteints ou non, résultats produits et leur qualité, valeur créée pour chaque partie prenante. Quantifiez autant que possible : ex. « La solution a réduit le temps de traitement de 40 \% » ou « Le dashboard a été adopté par 3 équipes en 2 semaines ». Assumez également les limites et ce qui n'a pas abouti.]}


\subsection{Analyse des impacts}

\begin{guideline}
Analysez l'impact de votre travail à différentes échelles : pour votre équipe directe, pour l'entreprise, et si pertinent pour des parties prenantes externes. Incluez également une réflexion sur la posture éthique que vous avez adoptée (respect des données, impacts environnementaux, inclusion...).
\end{guideline}

\instruction{[Analysez l'impact de vos productions à différentes échelles. Ex. : impact immédiat sur l'équipe, impact sur le produit ou le service, impact potentiel pour les utilisateurs finaux. Réfléchissez également aux implications éthiques de votre travail : confidentialité des données, accessibilité, empreinte carbone, diversité... Comment avez-vous intégré ces dimensions dans vos choix ?]}


\subsection{Analyse réflexive sur les compétences mobilisées}

\begin{guideline}
C'est la pièce maîtresse du dossier. Identifiez les compétences du référentiel que vous avez mobilisées, expliquez comment et dans quelle proportion. Croisez avec vos traces en annexe. Réfléchissez aussi aux compétences non mobilisées et pourquoi. Soyez spécifique : évitez les affirmations sans preuves.
\end{guideline}

\instruction{[Pour chaque compétence mobilisée, expliquez : dans quel contexte vous l'avez utilisée, comment vous l'avez mise en œuvre, quelles composantes essentielles ont été activées, et quelle preuve (trace) vous y associez. Ex. : « Dans le cadre de la refonte du module X, j'ai mobilisé la compétence [C2 --- Concevoir des solutions] en définissant des métriques d'évaluation comparatives (Annexe B.3). » Pour les compétences non mobilisées, justifiez pourquoi (contexte du stage ne s'y prêtait pas, périmètre limité...) et donnez une base réflexive de comment vous auriez pu le mettre en place.]}


\subsection{Projection professionnelle}

\begin{guideline}
Réfléchissez à votre futur rôle en tant qu'ingénieur·e. Qu'est-ce que ce stage a confirmé ou remis en question dans vos aspirations ? Quelles compétences souhaitez-vous développer en priorité ? Quel type d'environnement de travail vous correspond ? Adoptez une posture de recul et de projection.
\end{guideline}

\instruction{[Réfléchissez à la suite : qu'avez-vous appris sur vous-même en tant que futur·e ingénieur·e ? Ce stage a-t-il confirmé ou modifié votre orientation professionnelle ? Quelles compétences souhaitez-vous approfondir ? Quel type de poste, d'entreprise ou de secteur envisagez-vous ? Cette partie est personnelle : assumez vos choix et vos doutes.]}
